\clearpage
\section{Quantum vocabulary, introduced from zero}\label{sec:vocabulary}
\subsection{A bit, a qubit, a gate, and a measurement}
A classical bit has value $0$ or $1$. A qubit has two distinguished states,
written $\ket{0}$ and $\ket{1}$. The vertical bar and angle bracket form a
\emph{ket}: notation for a state vector. They do not mean absolute value here.
In column-vector notation,
\begin{equation}
\ket{0}=\begin{pmatrix}1\\0\end{pmatrix},\qquad
\ket{1}=\begin{pmatrix}0\\1\end{pmatrix}.
\end{equation}
A general pure one-qubit state is
\begin{equation}\label{eq:one-qubit-state}
\ket{\psi}=a\ket{0}+b\ket{1},\qquad |a|^2+|b|^2=1.
\end{equation}
The Greek letter $\psi$ (``psi'') names the state. The complex numbers $a,b$ are
\emph{amplitudes}. Measurement in this basis produces ordinary bit $0$ with
probability $|a|^2$, or bit $1$ with probability $|b|^2$. The squared magnitude
is why an amplitude is not itself a probability.

A \emph{gate} is an allowed transformation of the state, represented by a
unitary matrix. A \emph{circuit} is a sequence of gates. Measurement converts
the final state into a classical outcome. A \emph{shot} is one preparation,
execution, and measurement of the circuit.

\subsection{Superposition and phase}
The state
\begin{equation}
\ket{+}=\frac{\ket{0}+\ket{1}}{\sqrt 2}
\end{equation}
has equal measurement probabilities. Its two amplitudes have magnitude
$1/\sqrt2$, whose square is $1/2$. The square root and normalization matter.

An amplitude can be written $a=r e^{\ii\phi}$, where $r\ge0$ is its magnitude,
$\phi$ is its phase angle in radians, and $\ii^2=-1$. Then $|a|^2=r^2$.
An arrow in a plane is a useful picture: length represents magnitude and
direction represents phase. Changing only direction does not change length.
This is a mathematical picture, not a physical arrow in the chip.

Relative phases matter when later operations combine amplitude contributions.
Contributions may reinforce or cancel; this is \emph{interference}. For example,
$(\ket0+\ket1)/\sqrt2$ and $(\ket0-\ket1)/\sqrt2$ give the same probabilities
if measured immediately, but a subsequent Hadamard gate distinguishes them.
Appendix~\ref{app:math} works this out explicitly.

\subsection{Many qubits and the meaning of a bit string}
With $n$ qubits, a state is
\begin{equation}\label{eq:n-state}
\ket{\psi}=\sum_{\boldsymbol{x}\in\{0,1\}^n}a_{\boldsymbol{x}}\ket{\boldsymbol{x}},
\qquad \sum_{\boldsymbol{x}}|a_{\boldsymbol{x}}|^2=1.
\end{equation}
Here $\ket{\boldsymbol{x}}$ denotes the basis state in which qubit $q_i$ stores
bit $x_i$. The logical vector is ordered $(x_0,\ldots,x_{n-1})$; Qiskit's
\emph{displayed strings} reverse that order, as explained later. There are
$2^n$ possible outcomes: $2^5=32$ in our direct five-qubit encoding. One shot
returns one outcome, not a readable list of all 32 amplitudes.

An \emph{entangled} state cannot be written as a product of independent
single-qubit states. The RZZ interactions in our cost circuit can generate
entanglement. Entanglement by itself does not certify an optimization benefit.

\subsection{Operators, eigenvalues, and two different meanings of H}
An \emph{operator} is a linear transformation of state vectors, often written
as a matrix. If an operator $\widehat A$ acts on a state $\ket u$ as
\begin{equation}
\widehat A\ket u=\lambda\ket u,
\end{equation}
then $\ket u$ is an \emph{eigenvector} and $\lambda$ its \emph{eigenvalue}.
The output points along the same vector, multiplied by a number. This gives a
precise meaning to saying that a cost operator ``associates a score'' with a
definite booking string.

We distinguish the Hadamard gate $H_{\mathrm{Had}}$ from the cost Hamiltonian
$\HC$. They share a conventional letter but are different objects. A
Hamiltonian is a Hermitian operator; a gate is unitary. We use the Hamiltonian
to construct gates rather than applying its generally nonunitary matrix as a gate.

\section{Checkpoint 3a: translate binary costs into an Ising operator}\label{sec:ising}
\subsection{Pauli Z and the identity operator}
The identity operator $\Id$ leaves every state unchanged. The Pauli Z operator is
\begin{equation}
\widehat Z=\begin{pmatrix}1&0\\0&-1\end{pmatrix},\qquad
\widehat Z\ket0=+\ket0,\qquad \widehat Z\ket1=-\ket1.
\end{equation}
Its eigenvalues on the two bit states are $+1$ and $-1$. Applying Z to a
superposition changes the sign of the $\ket1$ amplitude; it does not flip a bit.
The operator $\widehat Z_i$ acts on qubit $i$ and as the identity on other qubits.

The classical sign $z_i$ associated with a definite bit is
\begin{equation}
z_i=(-1)^{x_i},\qquad x_i=\frac{1-z_i}{2}.
\end{equation}
\begin{center}
\begin{tabular}{ccc}
\toprule
Booking bit $x_i$ & Sign $z_i$ & $(1-z_i)/2$\\
\midrule
0 & $+1$ & 0\\
1 & $-1$ & 1\\
\bottomrule
\end{tabular}
\end{center}
The quantum counterpart of the binary variable is the projector
\begin{equation}\label{eq:projector}
\widehat x_i=\frac{\Id-\widehat Z_i}{2}.
\end{equation}
It has eigenvalue zero on bit $0$ and one on bit $1$. Keeping $x_i$, $z_i$, and
$\widehat Z_i$ distinct prevents a common confusion: a bit value is not a matrix.

\subsection{One cost term, step by step}
Consider only $f(x)=20x$. It assigns scores zero and twenty. Substituting gives
\begin{equation}
20x=10-10z,\qquad \HF=10\Id-10\widehat Z.
\end{equation}
Thus $\HF\ket0=0\ket0$ and $\HF\ket1=20\ket1$. The scoring function is
unchanged; the same scores are now eigenvalues of an operator. In this context,
``energy'' means optimization score, not the EV's battery energy in kWh.

\subsection{Translate every linear and pair term}
For the expanded QUBO in Equation~\eqref{eq:expanded-qubo},
\begin{align}
L_i x_i &\longmapsto \frac{L_i}{2}\Id-\frac{L_i}{2}\widehat Z_i,\\
Q_{ij}x_i x_j &\longmapsto \frac{Q_{ij}}4
\left(\Id-\widehat Z_i-\widehat Z_j+\widehat Z_i\widehat Z_j\right).
\end{align}
All terms can be collected as
\begin{equation}\label{eq:full-hamiltonian}
\HF=C_I\Id+\HC,\qquad
\HC=\sum_i h_i\widehat Z_i+\sum_{i<j}J_{ij}\widehat Z_i\widehat Z_j,
\end{equation}
with coefficients
\begin{align}
C_I&=C_0+\frac12\sum_iL_i+\frac14\sum_{i<j}Q_{ij},\label{eq:offset}\\
h_i&=-\frac{L_i}{2}-\frac14\sum_{j\ne i}Q_{\{i,j\}},\label{eq:h}\\
J_{ij}&=\frac{Q_{ij}}4.\label{eq:j}
\end{align}
In Equation~\eqref{eq:h}, $Q_{\{i,j\}}$ means the stored coefficient for the
unordered pair, using its smaller index first. An absent pair contributes zero.
This is precisely the bookkeeping in \code{qubo_to_ising()}.

The \emph{Ising form} uses signs $\pm1$ and their pair products. The name comes
from a model of interacting spins in physics. We use its mathematical structure;
no knowledge of physical magnets is required to formulate our application.

\subsection{The sample's actual operator}
The calculated single-qubit coefficients are
\begin{equation}
(h_0,h_1,h_2,h_3,h_4)
=(-83.0875,-70.975,-74.15,-76.15,4.5875).
\end{equation}
Every pair within $\{0,1,2,3\}$ has $J_{ij}=28.825$, and $J_{04}=14.4125$.
All other pair coefficients vanish. The offset is $C_I=227.7125$.
\begin{equation}\label{eq:actual-ising}
\begin{aligned}
\HC={}&-83.0875\widehat Z_0-70.975\widehat Z_1-74.15\widehat Z_2\\
&-76.15\widehat Z_3+4.5875\widehat Z_4\\
&+28.825\!\!\sum_{0\le i<j\le3}\widehat Z_i\widehat Z_j
+14.4125\widehat Z_0\widehat Z_4.
\end{aligned}
\end{equation}
Appendix~\ref{app:numeric} verifies these numbers against the original QUBO.
The Qiskit class \code{SparsePauliOp} stores this weighted sum without constructing
a full dense $32\times32$ matrix. ``Sparse'' means representing the listed terms;
it does not mean that the simulator's full statevector stops growing with $2^n$.

\subsection{Offset and normalization}
The circuit uses the normalized nonconstant cost
\begin{equation}\label{eq:normalized}
s=\max\bigl(\{|h_i|\}_i\cup\{|J_{ij}|\}_{i<j}\bigr),\qquad
\widetilde H_C=\frac{\HC}{s}.
\end{equation}
For this sample, $s=83.0875$. The implementation uses $s=1$ if all coefficients
vanish. Positive scaling and adding constants preserve the ordering of scores.
The reported expectation is restored by
\begin{equation}\label{eq:restore-score}
\mathbb E[f]=s\,\expect{\widetilde H_C}+C_I.
\end{equation}
The optimizer receives $\expect{\widetilde H_C}$; the printed history uses
Equation~\eqref{eq:restore-score}. These are different numerical scales for the
same optimization landscape.

\section{Checkpoint 3b: build the actual QAOA gates}\label{sec:gates}
\subsection{Initialize with Hadamard gates}
The single-qubit Hadamard gate is
\begin{equation}
H_{\mathrm{Had}}=\frac1{\sqrt2}\begin{pmatrix}1&1\\1&-1\end{pmatrix},
\qquad H_{\mathrm{Had}}\ket0=\ket+.
\end{equation}
Starting with all qubits in zero, apply one H gate per qubit:
\begin{equation}
\ket{\psi_0}=\ket+^{\otimes n}
=\frac1{\sqrt{2^n}}\sum_{\boldsymbol{x}}\ket{\boldsymbol{x}}.
\end{equation}
The tensor product $\otimes$ combines systems. Initially each five-bit outcome
has probability $1/32$. Initialization does not know which schedule is good.
\begin{lstlisting}[language=Python]
circuit.h(range(circuit.num_qubits))
\end{lstlisting}

\subsection{Cost operations change phases}
For one layer the cost operation is
\begin{equation}\label{eq:cost-unitary}
U_C(\gamma)=\exp\!\left(-\ii\gamma\widetilde H_C\right).
\end{equation}
The real number $\gamma$ (``gamma'') is adjustable. If
$\HC\ket{\boldsymbol{x}}=E_C(\boldsymbol{x})\ket{\boldsymbol{x}}$, then
\begin{equation}
U_C(\gamma)\ket{\boldsymbol{x}}
=e^{-\ii\gamma E_C(\boldsymbol{x})/s}\ket{\boldsymbol{x}}.
\end{equation}
The phase factor has magnitude one. It changes the amplitude's direction, not
its magnitude, so the cost layer alone leaves computational-basis probabilities
unchanged. A higher score does not automatically become a lower probability.

All our Z and ZZ terms commute, meaning their operator products can be reordered
without changing the result. Consequently Equation~\eqref{eq:cost-unitary}
factors exactly into one gate per term; no approximation from splitting these
commuting terms is required.

\subsection{Why RZ and RZZ have these angles}
Qiskit's rotation conventions are\cite{rz,rzz}
\begin{align}
\RZ(\theta)&=e^{-\ii\theta\widehat Z/2},\\
\RZZ(\theta)&=e^{-\ii\theta\widehat Z_i\widehat Z_j/2}.
\end{align}
The symbol $\theta$ (``theta'') is an angle in radians. One full turn is $2\pi$
radians; $\pi/2$ radians is $90^\circ$. Matching the exponents gives
\begin{align}
e^{-\ii\gamma h_i\widehat Z_i/s}
&=R_{Z,i}\!\left(\frac{2\gamma h_i}{s}\right),\\
e^{-\ii\gamma J_{ij}\widehat Z_i\widehat Z_j/s}
&=R_{ZZ,ij}\!\left(\frac{2\gamma J_{ij}}s\right).
\end{align}
This is the source of the factor of two in the code. There is no new vehicle
rule hidden in the gate angle; it is a convention of the gate definition.
\begin{lstlisting}[language=Python]
circuit.rz(2 * gamma[layer] * coefficient / scale, qubit)
circuit.rzz(2 * gamma[layer] * coefficient / scale, i, j)
\end{lstlisting}
For example, $h_0=-s$, so the first qubit's RZ angle is simply $-2\gamma$.

\subsection{Mixing turns phase differences into probability changes}
Pauli X flips the two basis states:
\begin{equation}
\widehat X=\begin{pmatrix}0&1\\1&0\end{pmatrix},\qquad
\widehat X\ket0=\ket1,\quad \widehat X\ket1=\ket0.
\end{equation}
The mixer Hamiltonian and its operation are
\begin{equation}
\widehat H_M=\sum_i\widehat X_i,\qquad
U_M(\beta)=e^{-\ii\beta\widehat H_M}
=\prod_i R_{X,i}(2\beta).
\end{equation}
The real number $\beta$ (``beta'') controls mixing. RX combines the amplitudes
for basis states differing in one bit. Depending on their phases, contributions
reinforce or cancel. Useful angles must be learned; arbitrary mixing can make
the distribution worse.
\begin{lstlisting}[language=Python]
for qubit in range(circuit.num_qubits):
    circuit.rx(2 * beta[layer], qubit)
\end{lstlisting}

\subsection{A one-qubit demonstration before five-qubit scheduling}
Consider again $f(x)=20x$. Ignore normalization for this small derivation, so
the nonconstant cost is $-10\widehat Z$. Choose $\gamma=\pi/40$ and
$\beta=-\pi/4$. The corresponding rotations are $\RZ(-\pi/2)$ and
$\RX(-\pi/2)$.
\begin{center}
\begin{tabular}{lcc}
\toprule
Circuit so far & $\Pr(0)$ & $\Pr(1)$\\
\midrule
$H_{\mathrm{Had}}$ & $1/2$ & $1/2$\\
$\RZ(-\pi/2)H_{\mathrm{Had}}$ & $1/2$ & $1/2$\\
$\RX(-\pi/2)\RZ(-\pi/2)H_{\mathrm{Had}}$ & $1$ & $0$\\
\bottomrule
\end{tabular}
\end{center}
These angles are selected for the demonstration. They are not a general solution
for the EV instance. Appendix~\ref{app:math} supplies the amplitude calculation.

\subsection{Layers, gate order, and the constant offset}
For $p$ layers, begin with $\ket{\psi_0}$ and repeat
\begin{equation}\label{eq:qaoa-recursion}
\ket{\psi_\ell}=U_M(\beta_\ell)U_C(\gamma_\ell)\ket{\psi_{\ell-1}},
\qquad \ell=1,\ldots,p.
\end{equation}
Operators on the right act first: cost, then mixing. With $p=2$ there are four
adjustable angles. Python indexes layers from zero, so mathematical
$\gamma_1$ is \code{gamma[0]}. The Hadamard initialization appears only once.
An \emph{ansatz} is this chosen parameterized circuit structure.

Including the offset $C_I\Id$ would multiply the whole state by
$e^{-\ii\gamma C_I/s}$, a common or \emph{global} phase. It cancels from all
measurement probabilities. Omitting it saves work without losing the cost
information needed to rank outcomes. We retain $C_I$ to restore reported scores.

\begin{checkpoint}
After H and the cost layer, why can all five-bit probabilities still be $1/32$?
Which operation allows phase differences to influence those probabilities?
\end{checkpoint}
